\paragraph{Proof of Theorem~\ref{thm:attack}.}
Let $\bar r=\frac1n\sum_t r_t$ and $\hat\Sigma$ the sample covariance. Then $Z_i=\sqrt n\,\bar r_i/\hat\sigma_i$, and by the CLT and $\hat\Sigma\to\Sigma$ a.s., $\sqrt n\,\bar r\Rightarrow\mathcal N(0,\Sigma)$, so $Z\Rightarrow Z^\ast\sim\mathcal N(0,\Sigma)$. Represent $Z^\ast=\sqrt\rho F\mathbf 1+\sqrt{1-\rho}\,\varepsilon$ with $F\sim\mathcal N(0,1)$ and $\varepsilon\sim\mathcal N(0,I_k)$ independent. Because $\operatorname{med}(Z^\ast)=\sqrt\rho F+\sqrt{1-\rho}\operatorname{med}(\varepsilon)$, the sign vector $s(Z^\ast)=\operatorname{sign}(\varepsilon-\operatorname{med}\varepsilon)$ is free of $F$. For even $k$ (ties have probability zero) exactly $k/2$ entries are $+1$, so $\mathbf 1^\top s=0$.

The reported portfolio $\hat w=s/k$ has $\sqrt n\,\hat w^\top\bar r=\frac1k s^\top(\sqrt n\,\bar r)$ and sample variance $\hat w^\top\hat\Sigma\hat w$. The map $x\mapsto s(x)$ is continuous at $Z^\ast$ almost surely, so by the continuous-mapping theorem and $\hat\Sigma\to\Sigma$,
\[
\tH(\hat w)=\frac{\frac1k s^\top\sqrt n\,\bar r}{\sqrt{\hat w^\top\hat\Sigma\hat w}}\;\Rightarrow\;\frac{\frac1k s^\top Z^\ast\cdot\sigma}{\sqrt{s^\top\Sigma s}/k},
\]
where we used $\sqrt n\,\bar r=\operatorname{diag}(\hat\sigma)Z$ with $\hat\sigma_i\to\sigma=1$. Now $s^\top\Sigma s=(1-\rho)\|s\|^2+\rho(\mathbf 1^\top s)^2=(1-\rho)k$ and $s^\top Z^\ast=\sqrt{1-\rho}\,s^\top\varepsilon=\sqrt{1-\rho}\sum_i|\varepsilon_i-\operatorname{med}\varepsilon|$, using $s^\top\mathbf 1\operatorname{med}\varepsilon=0$. Hence $\tH(\hat w)\Rightarrow k^{-1/2}\sum_i|\varepsilon_i-\operatorname{med}\varepsilon|$. Finally $\operatorname{med}\varepsilon\to0$ a.s.\ and $\frac1k\sum_i|\varepsilon_i-m|$ is $1$-Lipschitz in $m$, so $\frac1k\sum_i|\varepsilon_i-\operatorname{med}\varepsilon|\to\E|\varepsilon_1|=\sqrt{2/\pi}$ a.s., which gives $\tH(\hat w)/\sqrt{2k/\pi}\to1$.

For $\rho=0$ and $s_i=\operatorname{sign}(Z_i)$ the same argument gives $s^\top\Sigma s=k$ and $s^\top Z^\ast=\sum|\varepsilon_i|$, with the same limit. For keep-winners, $\hat w=\mathbf{1}\{Z>0\}/|S|$ with $|S|/k\to1/2$, mean $\sum_{i\in S}\varepsilon_i/|S|$ and variance $1/|S|$. The limit is $|S|^{-1/2}\sum_{i\in S}\varepsilon_i\approx\sqrt{k/2}\,\E[\varepsilon\mid\varepsilon>0]=\sqrt{k/2}\sqrt{2/\pi}=\sqrt{k/\pi}$. \qed

\paragraph{Proof of Corollary~\ref{cor:dsr}.}
The ledger's per-period Sharpe estimates are $Z_i/\sqrt n$ for $i\le k$ together with $\tH(\hat w)/\sqrt n$. Their sample variance satisfies $nV=\frac1{k}\big[\sum_i(Z_i-\bar Z)^2+(\tH(\hat w)-\bar Z)^2\big](1+o(1))$. The first term tends to $(1-\rho)$ (the cross-sectional variance of $Z^\ast$ given $F$) and the second to $(2k/\pi)/k=2/\pi$. With Gaussian returns $\hat\gamma_3\to0$ and $\hat\gamma_4\to3$, so the denominator of~\eqref{eq:dsr} tends to $\sqrt{1+\SR^2/2}=1+O(k/n)$. The numerator is $\sqrt{n-1}(\SR-\SR_0)=\tH(\hat w)-\sqrt{nV}\,\E\max_{k+1}+o(1)$, and $\E\max_{k+1}\le\sqrt{2\ln(k+1)}$. Therefore the DSR $z$-score is at least $\sqrt{2k/\pi}(1+o_p(1))-\sqrt{1-\rho+2/\pi}\sqrt{2\ln(k+1)}\to\infty$ in probability. \qed

\paragraph{Proof of Proposition~\ref{prop:tight}.}
A channel revealing $b$ bits has $|\cT|\le2^b$. By Theorem~\ref{thm:transcript} and the Gaussian tail $\bar\Phi^{-1}(p)=\sqrt{2\ln(1/p)}(1+o(1))$ as $p\to0$, the bar is $\sqrt{2(b\ln2+\ln(1/\alpha))}(1+o(1))=\sqrt{2b\ln 2}(1+o(1))$. The sign-flip analyst's report depends on $H$ only through $(\operatorname{sign}Z_i)_{i\le k}$, which is $k$ bits, and by Theorem~\ref{thm:attack} its $t$ is $\sqrt{2k/\pi}(1+o_p(1))$. Hence, for $b=o(n)$ and in the sequential limit of Theorem~\ref{thm:attack}, the supremum over $b$-bit analysts of the null $(1-\alpha)$-quantile of $\tH(\hat w)$ lies between $0.80\sqrt b$ and $1.18\sqrt b$. \qed

\paragraph{Proof of Proposition~\ref{prop:span}.} For any $w$ in the span, Cauchy--Schwarz in the $\hat\Sigma$ inner product gives $n(w^\top\bar r)^2/(w^\top\hat\Sigma w)\le n\,\bar r^\top\hat\Sigma^{-1}\bar r$. Under Gaussian $H_0$ the right-hand side is Hotelling's $T^2$, distributed as $\frac{k(n-1)}{n-k}F_{k,n-k}$ \citep{scheffe1953method,gibbons1989test}. The bound holds simultaneously for every $w$, hence for any data-dependent report. \qed

\paragraph{Remark (serial dependence and non-Gaussian rows).} With stationary, weakly dependent returns, Theorem~\ref{thm:attack} holds with $\Sigma$ replaced by the long-run covariance and a HAC $t$-statistic. Theorem~\ref{thm:transcript} needs the tail of the null $t$-statistic to be controlled at level $\alpha/|\cT|$. For i.i.d.\ non-Gaussian rows this follows from self-normalised moderate-deviation results when $c_\cT=o(n^{1/6})$ and the third moment is finite. For HAC statistics tail accuracy is poorer, and a block-bootstrap calibration of $\bar G$ is advisable when $|\cT|$ is large.
